\documentclass[12pt]{article}

\usepackage[a4paper,margin=1in]{geometry}
\usepackage{amsmath,amssymb}
\usepackage{setspace}
\usepackage{enumitem}

\setstretch{1.2}

\begin{document}

\begin{center}
\textbf{\Large CSE400 -- Fundamentals of Probability in Computing} \\[0.2cm]
\textbf{\large Lecture 7: Expectation, CDFs, PDFs and Problem Solving} \\[0.2cm]
\textbf{Instructor:} Dhaval Patel, PhD \\
\textbf{Date:} January 27, 2025
\end{center}

\vspace{0.5cm}

\section*{1. Cumulative Distribution Function (CDF)}

\subsection*{1.1 Definition of CDF}

Let $(X)$ be a random variable.

\[
F_X(x) = \Pr(X \le x), \quad -\infty < x < \infty
\]

\begin{itemize}
\item $(F_X(x))$ denotes the \textbf{Cumulative Distribution Function (CDF)} of random variable $(X)$.
\item Most of the information about the random experiment described by the random variable $(X)$ is determined by the behavior of $(F_X(x))$.
\end{itemize}

\subsection*{1.2 Properties of CDF}

For a valid CDF $(F_X(x))$, the following properties must hold:

\begin{enumerate}[label=\arabic*.]
\item \textbf{Boundedness}
\[
0 \le F_X(x) \le 1
\]

\item \textbf{Limits at infinity}
\[
F_X(-\infty) = 0
\]
\[
F_X(\infty) = 1
\]

\item \textbf{Monotonicity}

For $(x_1 < x_2)$,
\[
F_X(x_1) \le F_X(x_2)
\]

(Non-decreasing / monotonic increasing)

\item \textbf{Probability over an interval}

For $(x_1 < x_2)$,
\[
\Pr(x_1 < X \le x_2) = F_X(x_2) - F_X(x_1)
\]
\end{enumerate}

\section*{2. CDF -- Example \#1}

\subsection*{Problem:}

\textbf{Find the valid CDF}

\subsubsection*{(a)}

\[
F_X(x) = \frac{1}{2} + \frac{1}{\pi}\tan^{-1}(x)
\]

\begin{itemize}
\item Satisfies:
\begin{itemize}
\item $(0 \le F_X(x) \le 1)$
\item $(F_X(-\infty) = 0)$
\item $(F_X(\infty) = 1)$
\item Monotonic increasing
\end{itemize}
\end{itemize}

\textbf{Valid CDF}

\subsubsection*{(b)}

\[
F_X(x) = \left[1 - e^{-x}\right]u(x)
\]

\begin{itemize}
\item $(u(x))$: Unit Step Function
\item Defined as:
\[
u(x) =
\begin{cases}
0, & x < 0 \\
1, & x \ge 0
\end{cases}
\]
\end{itemize}

\textbf{Valid CDF}

\subsubsection*{(c)}

\[
F_X(x) = e^{-x^2}
\]

\textbf{Invalid CDF}

(Does not satisfy required CDF properties)

\subsubsection*{(d)}

\[
F_X(x) = x^2 u(x)
\]

\textbf{Invalid CDF}

\section*{3. CDF -- Example \#2}

\subsection*{Given:}

\[
F_X(x) = (1 - e^{-x})u(x)
\]

\subsection*{Find:}

\subsubsection*{(a) $\Pr(X > 5)$}

\[
\Pr(X > 5) = 1 - \Pr(X \le 5)
\]
\[
= 1 - F_X(5)
\]
\[
= e^{-5}
\]

\subsubsection*{(b) $\Pr(X < 5)$}

\[
\Pr(X < 5) = F_X(5)
\]

\subsubsection*{(c) $\Pr(3 < X < 7)$}

\[
\Pr(3 < X < 7) = F_X(7) - F_X(3)
\]

\subsubsection*{(d) $\Pr(X > 5 \mid X < 7)$}

Let:

\begin{itemize}
\item $(A = \{X > 5\})$
\item $(B = \{X < 7\})$
\end{itemize}

\[
\Pr(A \mid B) = \frac{\Pr(A \cap B)}{\Pr(B)}
\]

\[
= \frac{\Pr(5 < X < 7)}{\Pr(X < 7)}
\]

\[
= \frac{F_X(7) - F_X(5)}{F_X(7)}
\]

\section*{4. Probability Density Function (PDF)}

\subsection*{4.1 Definition of PDF}

For a \textbf{continuous random variable}, the PDF of $(X)$ evaluated at point $(x)$ is:

\[
f_X(x) = \lim_{\epsilon \to 0} \frac{\Pr(x \le X < x + \epsilon)}{\epsilon}
\]

\subsection*{4.2 PDF--CDF Relationship (Derivation)}

Recall for a continuous range:

\[
\Pr(x \le X < x + \epsilon) = F_X(x + \epsilon) - F_X(x)
\]

Substitute into PDF definition:

\[
f_X(x) = \lim_{\epsilon \to 0} \frac{F_X(x + \epsilon) - F_X(x)}{\epsilon}
\]

Taking the limit:

\[
\boxed{f_X(x) = \frac{dF_X(x)}{dx}}
\]

\subsection*{4.3 Important Result}

\begin{itemize}
\item The \textbf{PDF of a random variable is the derivative of its CDF}
\item Conversely, the \textbf{CDF of a random variable is the integral of its PDF}
\end{itemize}

\[
F_X(x) = \int_{-\infty}^{x} f_X(t)\,dt
\]

\section*{5. Relationship Between PDF and CDF}

\begin{itemize}
\item Area under the PDF curve from $(-\infty)$ to $(x)$ gives $(F_X(x))$
\item Graphical relationship shown in the slide:
\begin{itemize}
\item Shaded area under $(f_X(x))$ up to $(x_0)$ equals $(F_X(x_0))$
\end{itemize}
\end{itemize}

\section*{6. Expectation of Random Variables}

(Only headings present in slides -- content continues in later slides)

\subsection*{Topics listed:}

\begin{itemize}
\item Definition and Example
\item Expectation of a Function of RV
\item Linear Operation with Expectation
\item $(n^{\text{th}})$ Moments and Central Moments of RVs:
\begin{itemize}
\item Variance
\item Skewness
\item Kurtosis
\end{itemize}
\end{itemize}

(Details appear in subsequent slides beyond the provided content)

\end{document}
